\documentclass[a4paper,12pt]{article}

% --------------------------------------------------------
% CODIFICAÇÃO E TIPOGRAFIA
% --------------------------------------------------------
\usepackage[utf8]{inputenc}       % Permite acentuação direta no código-fonte
\usepackage[T1]{fontenc}          % Usa codificação T1 (melhor para português)
\usepackage[brazil]{babel}        % Tradução e hifenização para português
\usepackage{microtype}            % Melhora o espaçamento entre letras/palavras
\usepackage{parskip}              % Espaço entre parágrafos em vez de recuo

% --------------------------------------------------------
% PACOTES MATEMÁTICOS
% --------------------------------------------------------
\usepackage{amsmath}              % Ambientes matemáticos como align, equation
%\usepackage{amssymb}              % Símbolos matemáticos adicionais
\usepackage{amsthm}               % Ambientes de teoremas
\usepackage{mathtools}            % Extensões do amsmath (ex: \coloneqq)
\usepackage{mathrsfs}             % Fonte caligráfica com \mathscr
\usepackage{bbm}                  % Indicador \mathbbm{1}
\usepackage{dsfont}               % Alternativa para conjuntos: \mathds
\usepackage{commath}              % Notações como \abs{}, \norm{}, \eval{}
\usepackage{pgfmath}             % Permite realizar cálculos matemáticos (somas, produtos, números aleatórios, etc.)
\usepackage{siunitx}              % sistema de unidades
\sisetup{group-separator = {\,}} % espaço fino

% --------------------------------------------------------
% FONTE MODERNA (TEXTO E MATEMÁTICA)
% --------------------------------------------------------
%\usepackage{newtxtext}            % Fonte do texto (estilo Times)
%\usepackage{newtxmath}            % Fonte matemática compatível com pacotes AMS
\usepackage[bitstream-charter]{mathdesign} % Fonte elegante e matemática harmoniosa

% --------------------------------------------------------
% AMBIENTES TEÓRICOS PERSONALIZADOS
% --------------------------------------------------------
\theoremstyle{definition}
\newtheorem{definicao}{Definição}       % Definição numerada por seção
\newtheorem{exemplo}{Exemplo}           % Exemplo numerado por seção

%\theoremstyle{plain}
\newtheorem{teorema}{Teorema}           % Teorema numerado por seção
\newtheorem{lema}[teorema]{Lema}                 % Lema com numeração conjunta
\newtheorem{proposicao}[teorema]{Proposição}     % Proposição idem
\newtheorem{corolario}[teorema]{Corolário}       % Corolário idem
\newtheorem{exercicio}[teorema]{Exercício}

%\theoremstyle{remark}
\newtheorem{observacao}[teorema]{Observação}     % Observação com mesmo contador

% --------------------------------------------------------
% TABELAS E FIGURAS
% --------------------------------------------------------
\usepackage{graphicx}             % Inclusão de imagens
\usepackage{float}                % Controle de posicionamento (ex: [H])
\usepackage{caption}              % Personalização de legendas
\usepackage{subcaption}          % Subfiguras com \begin{subfigure}
\usepackage{booktabs}            % Tabelas com qualidade profissional
\usepackage{array}               % Mais opções em colunas de tabelas
\usepackage{multirow}            % Células que ocupam várias linhas
\usepackage{multicol}            % Disposição em colunas múltiplas
\usepackage{tabularx}            % Tabelas com colunas de largura ajustável
\usepackage{longtable}           % Tabelas que quebram página
\usepackage{arydshln}            % Linhas tracejadas horizontais e verticais
\usepackage{makecell}            % Células com múltiplas linhas (e quebra de linha)
\usepackage{diagbox}             % Cabeçalho diagonal em tabelas
\usepackage{pdflscape}           % Gira páginas (modo paisagem)
\newcolumntype{L}[1]{>{\raggedright\let\newline\\\arraybackslash\hspace{0pt}}m{#1}}
\newcolumntype{C}[1]{>{\centering\let\newline\\\arraybackslash\hspace{0pt}}m{#1}}
\newcolumntype{R}[1]{>{\raggedleft\let\newline\\\arraybackslash\hspace{0pt}}m{#1}}

% --------------------------------------------------------
% FORMATAÇÃO E ESTRUTURA
% --------------------------------------------------------
\usepackage{geometry}            % Configuração de margens
\geometry{a4paper, left=1cm, right=1cm, bottom=1.5cm, top=1.5cm, headsep=1cm, footskip=1cm}
\usepackage{titlesec}            % Controle sobre títulos de seções
\usepackage{fancyhdr}            % Cabeçalhos e rodapés personalizados
\pagestyle{fancy}
\fancyhf{}                % Limpa cabeçalho e rodapé
\fancyfoot[R]{\thepage}   % Numeração à direita no rodapé
\renewcommand{\headrulewidth}{0pt} % Remove linha no topo
\usepackage[shortlabels]{enumitem}            % Listas com controle de espaçamento e símbolo
\usepackage{tikz}                % Desenho de gráficos vetoriais (diagramas, etc.)
\usetikzlibrary{matrix,arrows,patterns,snakes,decorations.pathreplacing,3d,arrows.meta,calc}
\usepackage{etoolbox}             % Permite lógica condicional e manipulação de comandos
\usepackage{background}           % Permite adicionar conteúdo fixo no fundo da página (ex: linhas, marcas d'água)
\SetBgContents{} % Remove conteúdo padrão ("DRAFT")
\definecolor{background-color}{RGB}{233 227 206}

% --------------------------------------------------------
% GRÁFICOS E PLOTAGENS
% --------------------------------------------------------
\usepackage{pgfplots}             % Criação de gráficos vetoriais em 2D e 3D
\pgfplotsset{compat=1.18}         % Compatibilidade com versão atual do pacote
\usepackage{currfile} % Permite saber qual é o caminho completo do arquivo atual,

% --------------------------------------------------------
% LINKS E CORES
% --------------------------------------------------------
\usepackage{xcolor}              % Definição de cores (texto, links, etc.)
\usepackage[hidelinks]{hyperref} % Links sem moldura colorida no PDF
\usepackage{url}                 % Quebra automática de URLs
% Dica: se quiser links coloridos, use:
% \usepackage[colorlinks=true, linkcolor=blue, citecolor=blue, urlcolor=blue]{hyperref}

% --------------------------------------------------------
% CAIXAS E DESTAQUES VISUAIS
% --------------------------------------------------------
\usepackage[most]{tcolorbox}      % Criação de caixas coloridas e personalizadas para destaques, avisos, exemplos, etc.

% --------------------------------------------------------
% CÓDIGOS E ALGORITMOS
% --------------------------------------------------------
\usepackage{listings}            % Inclusão de código-fonte com destaque
\usepackage{algorithm}           % Ambiente de algoritmo
\usepackage[noend]{algpseudocode} % Sintaxe tipo pseudocódigo

% --------------------------------------------------------
% COMANDOS
% --------------------------------------------------------

\newcommand{\bbN}{\mathbb{N}}
\newcommand{\bbZ}{\mathbb{Z}}
\newcommand{\bbQ}{\mathbb{Q}}
\newcommand{\bbR}{\mathbb{R}}
\newcommand{\calnN}{\mathcal{N}}

\newcommand{\assunto}{ESTATÍSTICA DESCRITIVA E PROBABILIDADE}
\newcommand{\atividade}{prova} % prova ou lista
\usepackage{xfp}
\newcommand{\ValorProva}{33}
\newcommand{\NumTotalQuestoes}{5}
\newcommand{\ValorQ}{%
    \frac{\inteval{\ValorProva / gcd(\ValorProva,\NumTotalQuestoes)}}
         {\inteval{\NumTotalQuestoes / gcd(\ValorProva,\NumTotalQuestoes)}}%
}
\newcommand{\turma}{ESTATÍSTICA}
\newcommand{\professor}{Igor Martins Silva}
\newcommand{\semestre}{2}
\newcommand{\ano}{2026}
\newcommand{\mesTexto}{outubro}
\newcommand{\mesNum}{10}
\newcommand{\dia}{14}
\newcommand{\dataTexto}{\dia\ de \mesTexto\ de \ano}
\newcommand{\dataNun}{\dia/\mesNum/\ano}

\ifdefstring{\atividade}{prova}{
    \newcommand{\titulo}{Avaliação de Estatística}
}{
    \newcommand{\titulo}{Lista de exercícios de Estatística}
}

\edef\pathCapa{\currfiledir}

\newcommand{\subtitulo}{
    %
    \noindent
    \begin{minipage}{2.8cm}
        \includegraphics[width=2.8cm]{\pathCapa Logo_Cefet.png}
    \end{minipage}
    \hfill
    \begin{minipage}{\dimexpr\textwidth-6.2cm\relax}
        \begin{center}
            CEFET-MG -- Campus Contagem \\[3pt]
            \titulo \ -- \semestre.\textordmasculine\ Semestre de \ano \\[3pt]
            \professor
        \end{center}
    \end{minipage}
    \hfill
    \begin{minipage}{3.2cm}
        \includegraphics[width=3.2cm]{\pathCapa Brasao_Brasil.png}
    \end{minipage}
    %
    
    %
    \begin{center}
        \vspace{15pt}
        \begin{tabular}{|C{250pt}|C{250pt}|}
            \hline
            \textbf{ASSUNTO} &
            \textbf{TURMA} \\
            \hline
            \assunto & \turma \\
            \hline
        \end{tabular}
    \end{center}
    %
    \vspace{15pt}
}

\newcommand{\subtituloAlt}{
    %
    \noindent
    \begin{minipage}{2.8cm}
        \includegraphics[width=2.8cm]{\pathCapa Logo_Cefet.png}
    \end{minipage}
    \hfill
    \begin{minipage}{\dimexpr\textwidth-6.2cm\relax}
        \begin{center}
            CEFET-MG -- Campus Contagem \\[3pt]
            \titulo \ -- \semestre.\textordmasculine\ Semestre de \ano \\[3pt]
            \professor
        \end{center}
    \end{minipage}
    \hfill
    \begin{minipage}{3.2cm}
        \includegraphics[width=3.2cm]{\pathCapa Brasao_Brasil.png}
    \end{minipage}
    %
    
    %
    \begin{center}
        \vspace{15pt}
        \begin{tabular}{|C{250pt}|C{100pt}|C{100pt}|}
            \hline
            \textbf{ASSUNTO} &
            \textbf{DATA} &
            \textbf{VALOR} \\
            \hline
            \assunto & \dataNun & \ValorProva\ pontos \\
            \hline
        \end{tabular}
    \end{center}
    %
    \vspace{15pt}
}

\newcommand{\valorquestao}{}

\newcommand{\definevalorquestao}{%
    \ifdefstring{\atividade}{lista}{%
        % Não faz nada para lista
    }{%
        \renewcommand{\valorquestao}{\ValorQ pontos}
    }%
}

\newenvironment{exercicioBanco}[1][]{%
    \ifdefstring{\atividade}{prova}{%
        \begin{exercicio}[#1]%
    }{%
        \begin{exercicio}%
    }%
}{%
    \end{exercicio}
}

\ifdefstring{\atividade}{lista}{
    \pagecolor{background-color}
}{
    
}



\hypersetup{
    pdfauthor={\professor},
    pdftitle={\titulo -- \assunto},
}

\title{\titulo -- \assunto}
\author{\professor}
\date{\data}

%************* INÍCIO DO DOCUMENTO *************
\begin{document}
        
    \subtituloAlt
    
    \phantomsection
    \addcontentsline{toc}{section}{Exercício 1}
    %%%%%%%%%%%%%%%%%%%%%%%%%%%%%%%%%
%
%       EXERCÍCIO
%
%%%%%%%%%%%%%%%%%%%%%%%%%%%%%%%%%

\definevalorquestao

\begin{exercicioBanco}[\valorquestao]
A tabela a seguir apresenta a distribuição de frequências relativas aos tempos de atendimento (em minutos) registrados em um posto de saúde:

\begin{center}
\begin{tabular}{|c|c|}
\hline
Tempo (min) & Frequência ($f_i$) \\ \hline
$0 \vdash 10$ & 3 \\ \hline
$10 \vdash 20$ & 5 \\ \hline
$20 \vdash 30$ & 8 \\ \hline
$30 \vdash 40$ & 4 \\ \hline
$40 \vdash 50$ & 2 \\ \hline
\end{tabular}
\end{center}

Com base nessa tabela, classifique os dados quanto sua assimetria, usando o 2.\textordmasculine\ coeficiente de assimetria de Pearson, \(e_2\). Para isso, considere a seguinte classificação:
\begin{center}
    \renewcommand{\arraystretch}{1.2}
    \begin{tabular}{|c|c|}
        \hline
        \textbf{Valor de } $e_2$ & \textbf{Classificação da Assimetria} \\ \hline
        $|e_2| = 0$ & Simétrica \\ \hline
        $0 < |e_2| \le 0{,}15$ & Assimetria fraca \\ \hline
        $0{,}15 < |e_2| \le 1$ & Assimetria moderada \\ \hline
        $|e_2| > 1$ & Assimetria forte \\ \hline
    \end{tabular}
\end{center}
\end{exercicioBanco}

    
    \phantomsection
    \addcontentsline{toc}{section}{Exercício 2}
    %%%%%%%%%%%%%%%%%%%%%%%%%%%%%%%%%
%
%       EXERCÍCIO
%
%%%%%%%%%%%%%%%%%%%%%%%%%%%%%%%%%

\ifdefstring{\atividade}{prova}{%
    \renewcommand{\valorquestao}{\ValorQ\ pontos}
}%

\begin{exercicioBanco}[\valorquestao]
O volume diário de leite produzido por uma cooperativa agrícola (em milhares de litros) foi registrado durante 40 dias consecutivos e é apresentado no rol a seguir:

\begin{center}
    \renewcommand{\arraystretch}{1.3}
    \begin{tabular}{|C{30pt}|C{30pt}|C{30pt}|C{30pt}|C{30pt}|C{30pt}|C{30pt}|C{30pt}|C{30pt}|C{30pt}|}
        \hline
        10.2 & 10.4 & 10.5 & 10.7 & 11.0 & 11.1 & 11.3 & 11.4 & 11.5 & 11.5 \\
        \hline
        11.6 & 11.6 & 11.7 & 11.7 & 11.8 & 11.8 & 11.8 & 11.9 & 11.9 & 12.0 \\
        \hline
        12.0 & 12.1 & 12.1 & 12.2 & 12.2 & 12.3 & 12.3 & 12.4 & 12.4 & 12.5 \\
        \hline
        12.6 & 12.7 & 12.8 & 13.0 & 13.1 & 13.3 & 13.5 & 13.6 & 13.8 & 14.2 \\
        \hline
    \end{tabular}
\end{center}

\begin{enumerate}[(a)]
    \item Construa o histograma dessa variável estatística, considerando intervalos de classe com amplitude igual a \(0.8\), iniciando a primeira classe em \(10.0\).
    %
    \item Construa o boxplot (diagrama de caixas) correspondente a esse conjunto de dados.
\end{enumerate}
%
\end{exercicioBanco}


    
    \phantomsection
    \addcontentsline{toc}{section}{Exercício 3}
    %%%%%%%%%%%%%%%%%%%%%%%%%%%%%%%%%
%
%       EXERCÍCIO
%
%%%%%%%%%%%%%%%%%%%%%%%%%%%%%%%%%

\definevalorquestao

\begin{exercicioBanco}[\valorquestao]
Sejam \(A\) e \(B\) eventos de um mesmo espaço amostral tais que \(P(A) = 0{,}6\), \(P(B) = 0{,}5\) e \(P(A \cap B) = 0{,}4\). Calcule:

\begin{multicols}{3}
    \begin{enumerate}[(a)]
        \item \(P(A \cup B)\);
        %
        \item \(P(A \mid B)\);
        %
        \item \(P(A^{\complement} \mid B)\).
    \end{enumerate}
\end{multicols}
%
\end{exercicioBanco}


    
    \phantomsection
    \addcontentsline{toc}{section}{Exercício 4}
    %%%%%%%%%%%%%%%%%%%%%%%%%%%%%%%%%
%
%       EXERCÍCIO
%
%%%%%%%%%%%%%%%%%%%%%%%%%%%%%%%%%

\definevalorquestao

\begin{exercicioBanco}[\valorquestao]
Um tetraedro regular e justo (com 4 faces equiprováveis) possui suas faces pintadas com as seguintes cores:
\begin{multicols}{2}
\begin{itemize}
    \item \textbf{Face 1:} vermelho;
    \item \textbf{Face 2:} azul;
    \item \textbf{Face 3:} verde;
    \item \textbf{Face 4:} dividida em três regiões, uma vermelha, uma azul e uma verde.
\end{itemize}
\end{multicols}

\noindent O tetraedro é lançado e a face em contato com a mesa é observada. Considere os seguintes eventos:
\begin{align*}
    V &: \text{``A face sorteada contém a cor Vermelha''} \\
    A &: \text{``A face sorteada contém a cor Azul''} \\
    G &: \text{``A face sorteada contém a cor Verde''}
\end{align*}

\begin{enumerate}
    \item[(a)] Determine as probabilidades individuais de cada evento: $P(V)$, $P(A)$ e $P(G)$.
    \item[(b)] Mostre que os eventos $V$, $A$ e $G$ são independentes dois a dois.
    \item[(c)] Os eventos $V$, $A$ e $G$ são independentes? Justifique.
\end{enumerate}
\end{exercicioBanco}


    
    \phantomsection
    \addcontentsline{toc}{section}{Exercício 5}
    %%%%%%%%%%%%%%%%%%%%%%%%%%%%%%%%%
%
%       EXERCÍCIO
%
%%%%%%%%%%%%%%%%%%%%%%%%%%%%%%%%%

\ifdefstring{\atividade}{prova}{%
    \renewcommand{\valorquestao}{\ValorQ\ pontos}
}%

\begin{exercicioBanco}[\valorquestao]
Uma indústria automobilística adquire peças de três fornecedores distintos: \(A\), \(B\) e \(C\). O fornecedor \(A\) é responsável por \(20\%\) das peças compradas, o fornecedor \(B\) por \(30\%\) e o fornecedor \(C\) por \(50\%\). Historicamente, a probabilidade de uma peça apresentar algum defeito de fabricação é de \(0{,}05\) para o fornecedor \(A\), \(0{,}03\) para o fornecedor \(B\) e \(0{,}01\) para o fornecedor \(C\).

\begin{enumerate}[(a)]
    \item Qual é a probabilidade de uma peça selecionada aleatoriamente do estoque não apresentar defeito?
    %
    \item Se uma peça escolhida ao acaso apresentou defeito, qual é a probabilidade de ela ter sido produzida pelo fornecedor \(A\)?
\end{enumerate}
%
\end{exercicioBanco}


    
    %\newpage 
    
    %%%%%%%%%%%%%%%%%%%%%%%%%%%%%%%%%%
%
%       NORMAL PADRÃO
%
%%%%%%%%%%%%%%%%%%%%%%%%%%%%%%%%%

\phantomsection
\addcontentsline{toc}{section}{Tabela da normal padrão}

%\captionsetup[table]{labelformat=empty} % remove "Tabela 1" do título

\begin{center}
\textbf{Distribuição Normal: valores de $p$ tais que $P(0 \le Z \le z_c) = p$}
\end{center}

\vspace{20pt}

\setlength{\tabcolsep}{6pt}
\renewcommand{\arraystretch}{1.2}
\begin{tabular}{cc|cccccccccc}
    & & \multicolumn{10}{c}{\textbf{Segunda decimal de $z_c$}} \\[5pt]
    \multirow{38}{*}{\rotatebox{90}{\textbf{Parte inteira e primeira decimal de $z_c$}}} 
    & & \textbf{0} & \textbf{1} & \textbf{2} & \textbf{3} & \textbf{4} & \textbf{5} & \textbf{6} & \textbf{7} & \textbf{8} & \textbf{9} \\
    \midrule
    %\endfirsthead
    %\multicolumn{12}{c}%
    %{{\tablename\ \thetable{} -- continuação}} \\
    %\toprule
    %& & \multicolumn{10}{c}{\textbf{Segunda decimal de $z_c$}} \\
    %\multirow{38}{*}{\rotatebox{90}{\textbf{Parte inteira e primeira decimal de $z_c$}}} 
    %& & \textbf{0} & \textbf{1} & \textbf{2} & \textbf{3} & \textbf{4} & \textbf{5} & \textbf{6} & \textbf{7} & \textbf{8} & \textbf{9} \\
    %\midrule
    %\endhead
    %\bottomrule
    %\endfoot

    & \textbf{0.0} & 0.0000 & 0.0040 & 0.0080 & 0.0120 & 0.0160 & 0.0199 & 0.0239 & 0.0279 & 0.0319 & 0.0359 \\
    & \textbf{0.1} & 0.0398 & 0.0438 & 0.0478 & 0.0517 & 0.0557 & 0.0596 & 0.0636 & 0.0675 & 0.0714 & 0.0753 \\
    & \textbf{0.2} & 0.0793 & 0.0832 & 0.0871 & 0.0910 & 0.0948 & 0.0987 & 0.1026 & 0.1064 & 0.1103 & 0.1141 \\
    & \textbf{0.3} & 0.1179 & 0.1217 & 0.1255 & 0.1293 & 0.1331 & 0.1368 & 0.1406 & 0.1443 & 0.1480 & 0.1517 \\
    & \textbf{0.4} & 0.1554 & 0.1591 & 0.1628 & 0.1664 & 0.1700 & 0.1736 & 0.1772 & 0.1808 & 0.1844 & 0.1879 \\
    & \textbf{0.5} & 0.1915 & 0.1950 & 0.1985 & 0.2019 & 0.2054 & 0.2088 & 0.2123 & 0.2157 & 0.2190 & 0.2224 \\
    & \textbf{0.6} & 0.2257 & 0.2291 & 0.2324 & 0.2357 & 0.2389 & 0.2422 & 0.2454 & 0.2486 & 0.2517 & 0.2549 \\
    & \textbf{0.7} & 0.2580 & 0.2611 & 0.2642 & 0.2673 & 0.2704 & 0.2734 & 0.2764 & 0.2794 & 0.2823 & 0.2852 \\
    & \textbf{0.8} & 0.2881 & 0.2910 & 0.2939 & 0.2967 & 0.2995 & 0.3023 & 0.3051 & 0.3078 & 0.3106 & 0.3133 \\
    & \textbf{0.9} & 0.3159 & 0.3186 & 0.3212 & 0.3238 & 0.3264 & 0.3289 & 0.3315 & 0.3340 & 0.3365 & 0.3389 \\
    & \textbf{1.0} & 0.3413 & 0.3438 & 0.3461 & 0.3485 & 0.3508 & 0.3531 & 0.3554 & 0.3577 & 0.3599 & 0.3621 \\
    & \textbf{1.1} & 0.3643 & 0.3665 & 0.3686 & 0.3708 & 0.3729 & 0.3749 & 0.3770 & 0.3790 & 0.3810 & 0.3830 \\
    & \textbf{1.2} & 0.3849 & 0.3869 & 0.3888 & 0.3907 & 0.3925 & 0.3944 & 0.3962 & 0.3980 & 0.3997 & 0.4015 \\
    & \textbf{1.3} & 0.4032 & 0.4049 & 0.4066 & 0.4082 & 0.4099 & 0.4115 & 0.4131 & 0.4147 & 0.4162 & 0.4177 \\
    & \textbf{1.4} & 0.4192 & 0.4207 & 0.4222 & 0.4236 & 0.4251 & 0.4265 & 0.4279 & 0.4292 & 0.4306 & 0.4319 \\
    & \textbf{1.5} & 0.4332 & 0.4345 & 0.4357 & 0.4370 & 0.4382 & 0.4394 & 0.4406 & 0.4418 & 0.4429 & 0.4441 \\
    & \textbf{1.6} & 0.4452 & 0.4463 & 0.4474 & 0.4484 & 0.4495 & 0.4505 & 0.4515 & 0.4525 & 0.4535 & 0.4545 \\
    & \textbf{1.7} & 0.4554 & 0.4564 & 0.4573 & 0.4582 & 0.4591 & 0.4599 & 0.4608 & 0.4616 & 0.4625 & 0.4633 \\
    & \textbf{1.8} & 0.4641 & 0.4649 & 0.4656 & 0.4664 & 0.4671 & 0.4678 & 0.4686 & 0.4693 & 0.4699 & 0.4706 \\
    & \textbf{1.9} & 0.4713 & 0.4719 & 0.4726 & 0.4732 & 0.4738 & 0.4744 & 0.4750 & 0.4756 & 0.4761 & 0.4767 \\
    & \textbf{2.0} & 0.4772 & 0.4778 & 0.4783 & 0.4788 & 0.4793 & 0.4798 & 0.4803 & 0.4808 & 0.4812 & 0.4817 \\
    & \textbf{2.1} & 0.4821 & 0.4826 & 0.4830 & 0.4834 & 0.4838 & 0.4842 & 0.4846 & 0.4850 & 0.4854 & 0.4857 \\
    & \textbf{2.2} & 0.4861 & 0.4864 & 0.4868 & 0.4871 & 0.4875 & 0.4878 & 0.4881 & 0.4884 & 0.4887 & 0.4890 \\
    & \textbf{2.3} & 0.4893 & 0.4896 & 0.4898 & 0.4901 & 0.4904 & 0.4906 & 0.4909 & 0.4911 & 0.4913 & 0.4916 \\
    & \textbf{2.4} & 0.4918 & 0.4920 & 0.4922 & 0.4925 & 0.4927 & 0.4929 & 0.4931 & 0.4932 & 0.4934 & 0.4936 \\
    & \textbf{2.5} & 0.4938 & 0.4940 & 0.4941 & 0.4943 & 0.4945 & 0.4946 & 0.4948 & 0.4949 & 0.4951 & 0.4952 \\
    & \textbf{2.6} & 0.4953 & 0.4955 & 0.4956 & 0.4957 & 0.4959 & 0.4960 & 0.4961 & 0.4962 & 0.4963 & 0.4964 \\
    & \textbf{2.7} & 0.4965 & 0.4966 & 0.4967 & 0.4968 & 0.4969 & 0.4970 & 0.4971 & 0.4972 & 0.4973 & 0.4974 \\
    & \textbf{2.8} & 0.4974 & 0.4975 & 0.4976 & 0.4977 & 0.4977 & 0.4978 & 0.4979 & 0.4979 & 0.4980 & 0.4981 \\
    & \textbf{2.9} & 0.4981 & 0.4982 & 0.4982 & 0.4983 & 0.4984 & 0.4984 & 0.4985 & 0.4985 & 0.4986 & 0.4986 \\
    & \textbf{3.0} & 0.4987 & 0.4987 & 0.4987 & 0.4988 & 0.4988 & 0.4989 & 0.4989 & 0.4989 & 0.4990 & 0.4990 \\
    & \textbf{3.1} & 0.4990 & 0.4991 & 0.4991 & 0.4991 & 0.4992 & 0.4992 & 0.4992 & 0.4992 & 0.4993 & 0.4993 \\
    & \textbf{3.2} & 0.4993 & 0.4993 & 0.4994 & 0.4994 & 0.4994 & 0.4994 & 0.4994 & 0.4995 & 0.4995 & 0.4995 \\
    & \textbf{3.3} & 0.4995 & 0.4995 & 0.4995 & 0.4996 & 0.4996 & 0.4996 & 0.4996 & 0.4996 & 0.4996 & 0.4997 \\
    & \textbf{3.4} & 0.4997 & 0.4997 & 0.4997 & 0.4997 & 0.4997 & 0.4997 & 0.4997 & 0.4997 & 0.4997 & 0.4998 \\
    & \textbf{3.5} & 0.4998 & 0.4998 & 0.4998 & 0.4998 & 0.4998 & 0.4998 & 0.4998 & 0.4998 & 0.4998 & 0.4998 \\
    \bottomrule
\end{tabular}

%
\begin{comment}
%
\vspace{13pt}

%
\textbf{Resolução}.
%

%
\hfill\(\boxtimes\)
\end{comment}

    
    %\newpage
    
    %%\documentclass[tikz, border=10pt]{standalone}
\phantomsection
\addcontentsline{toc}{section}{Tabela da t-Student}

\captionsetup[table]{labelformat=empty} % remove "Tabela 1" do título

\begin{center}
\textbf{Distribuição \(t\)-Student: Valores de \(t_{c}\) tais que \(P(-t_{c} \le T \le t_{c}) = 1 - p\)}
\end{center}
\vspace{1em}
\renewcommand{\arraystretch}{1.1} % aumenta altura em 50%
%\footnotesize
\resizebox{\textwidth}{!}{%
\begin{tabular}{cc|ccccccccccccccc}
    & & \multicolumn{15}{c}{\textbf{Valor de \(p\)}} \\[10pt]
    \multirow{36}{*}{\rotatebox{90}{\textbf{Grau de liberdade}}}
    & & \textbf{90\%} & \textbf{80\%} & \textbf{70\%} & \textbf{60\%} & \textbf{50\%} & \textbf{40\%} & \textbf{30\%} & \textbf{20\%} & \textbf{10\%} & \textbf{5\%} & \textbf{4\%} & \textbf{2\%} & \textbf{1\%} & \textbf{0.2\%} & \textbf{0.1\%} \\
    \hline
    & \textbf{1} & 0.158 & 0.365 & 0.510 & 0.727 & 1.000 & 1.376 & 1.963 & 3.078 & 6.314 & 12.706 & 15.894 & 31.821 & 63.656 & 318.289 & 636.578 \\
    & \textbf{2} & 0.142 & 0.289 & 0.445 & 0.617 & 0.816 & 1.061 & 1.386 & 1.886 & 2.920 & 4.303 & 4.849 & 6.965 & 9.925 & 22.328 & 31.600 \\
    & \textbf{3} & 0.137 & 0.277 & 0.424 & 0.584 & 0.765 & 0.978 & 1.250 & 1.638 & 2.353 & 3.182 & 3.482 & 4.541 & 5.841 & 10.214 & 12.924 \\
    & \textbf{4} & 0.134 & 0.271 & 0.414 & 0.569 & 0.741 & 0.941 & 1.190 & 1.533 & 2.132 & 2.776 & 2.999 & 3.747 & 4.604 & 7.173 & 8.610 \\
    & \textbf{5} & 0.132 & 0.267 & 0.408 & 0.559 & 0.727 & 0.920 & 1.156 & 1.476 & 2.015 & 2.571 & 2.757 & 3.365 & 4.032 & 5.894 & 6.869 \\
    & \textbf{6} & 0.131 & 0.265 & 0.404 & 0.553 & 0.718 & 0.906 & 1.134 & 1.440 & 1.943 & 2.447 & 2.612 & 3.143 & 3.707 & 5.208 & 5.959 \\
    & \textbf{7} & 0.130 & 0.263 & 0.402 & 0.549 & 0.711 & 0.896 & 1.119 & 1.415 & 1.895 & 2.365 & 2.517 & 2.998 & 3.499 & 4.785 & 5.408 \\
    & \textbf{8} & 0.130 & 0.262 & 0.399 & 0.546 & 0.706 & 0.889 & 1.108 & 1.397 & 1.860 & 2.306 & 2.449 & 2.896 & 3.355 & 4.501 & 5.041 \\
    & \textbf{9} & 0.129 & 0.261 & 0.398 & 0.543 & 0.703 & 0.883 & 1.100 & 1.383 & 1.833 & 2.262 & 2.398 & 2.821 & 3.250 & 4.297 & 4.781 \\
    & \textbf{10} & 0.129 & 0.260 & 0.397 & 0.542 & 0.700 & 0.879 & 1.093 & 1.372 & 1.812 & 2.228 & 2.359 & 2.764 & 3.169 & 4.144 & 4.587 \\
    & \textbf{11} & 0.129 & 0.260 & 0.396 & 0.540 & 0.697 & 0.876 & 1.088 & 1.363 & 1.796 & 2.201 & 2.328 & 2.718 & 3.106 & 4.025 & 4.437 \\
    & \textbf{12} & 0.128 & 0.259 & 0.395 & 0.539 & 0.695 & 0.873 & 1.083 & 1.356 & 1.782 & 2.179 & 2.303 & 2.681 & 3.055 & 3.930 & 4.318 \\
    & \textbf{13} & 0.128 & 0.259 & 0.394 & 0.538 & 0.694 & 0.870 & 1.079 & 1.350 & 1.771 & 2.160 & 2.282 & 2.650 & 3.012 & 3.852 & 4.221 \\
    & \textbf{14} & 0.128 & 0.258 & 0.393 & 0.537 & 0.692 & 0.868 & 1.076 & 1.345 & 1.761 & 2.145 & 2.264 & 2.624 & 2.977 & 3.787 & 4.140 \\
    & \textbf{15} & 0.128 & 0.258 & 0.393 & 0.536 & 0.691 & 0.866 & 1.074 & 1.341 & 1.753 & 2.131 & 2.249 & 2.602 & 2.947 & 3.733 & 4.073 \\
    & \textbf{16} & 0.128 & 0.258 & 0.392 & 0.535 & 0.690 & 0.865 & 1.071 & 1.337 & 1.746 & 2.120 & 2.235 & 2.583 & 2.921 & 3.686 & 4.015 \\
    & \textbf{17} & 0.128 & 0.257 & 0.392 & 0.534 & 0.689 & 0.863 & 1.069 & 1.333 & 1.740 & 2.110 & 2.224 & 2.567 & 2.898 & 3.646 & 3.965 \\
    & \textbf{18} & 0.127 & 0.257 & 0.392 & 0.534 & 0.688 & 0.862 & 1.067 & 1.330 & 1.734 & 2.101 & 2.214 & 2.552 & 2.878 & 3.610 & 3.922 \\
    & \textbf{19} & 0.127 & 0.257 & 0.391 & 0.533 & 0.688 & 0.861 & 1.066 & 1.328 & 1.729 & 2.093 & 2.205 & 2.539 & 2.861 & 3.579 & 3.883 \\
    & \textbf{20} & 0.127 & 0.257 & 0.391 & 0.533 & 0.687 & 0.860 & 1.064 & 1.325 & 1.725 & 2.086 & 2.197 & 2.528 & 2.845 & 3.552 & 3.850 \\
    & \textbf{21} & 0.127 & 0.257 & 0.391 & 0.532 & 0.686 & 0.859 & 1.063 & 1.323 & 1.721 & 2.080 & 2.189 & 2.518 & 2.831 & 3.527 & 3.819 \\
    & \textbf{22} & 0.127 & 0.256 & 0.390 & 0.532 & 0.686 & 0.858 & 1.061 & 1.321 & 1.717 & 2.074 & 2.183 & 2.508 & 2.819 & 3.505 & 3.792 \\
    & \textbf{23} & 0.127 & 0.256 & 0.390 & 0.532 & 0.685 & 0.858 & 1.060 & 1.319 & 1.714 & 2.069 & 2.177 & 2.500 & 2.807 & 3.485 & 3.768 \\
    & \textbf{24} & 0.127 & 0.256 & 0.390 & 0.531 & 0.685 & 0.857 & 1.059 & 1.318 & 1.711 & 2.064 & 2.172 & 2.492 & 2.797 & 3.467 & 3.745 \\
    & \textbf{25} & 0.127 & 0.256 & 0.390 & 0.531 & 0.684 & 0.856 & 1.058 & 1.316 & 1.708 & 2.060 & 2.167 & 2.485 & 2.787 & 3.450 & 3.725 \\
    & \textbf{26} & 0.127 & 0.256 & 0.390 & 0.531 & 0.684 & 0.856 & 1.058 & 1.315 & 1.706 & 2.056 & 2.162 & 2.479 & 2.779 & 3.435 & 3.707 \\
    & \textbf{27} & 0.127 & 0.256 & 0.389 & 0.531 & 0.684 & 0.855 & 1.057 & 1.314 & 1.703 & 2.052 & 2.158 & 2.473 & 2.771 & 3.421 & 3.689 \\
    & \textbf{28} & 0.127 & 0.256 & 0.389 & 0.530 & 0.683 & 0.855 & 1.056 & 1.313 & 1.701 & 2.048 & 2.154 & 2.467 & 2.763 & 3.408 & 3.674 \\
    & \textbf{29} & 0.127 & 0.256 & 0.389 & 0.530 & 0.683 & 0.854 & 1.055 & 1.311 & 1.699 & 2.045 & 2.150 & 2.462 & 2.756 & 3.396 & 3.660 \\
    & \textbf{30} & 0.127 & 0.256 & 0.389 & 0.530 & 0.683 & 0.854 & 1.055 & 1.310 & 1.697 & 2.042 & 2.147 & 2.457 & 2.750 & 3.385 & 3.646 \\
    & \textbf{35} & 0.127 & 0.255 & 0.388 & 0.529 & 0.682 & 0.852 & 1.052 & 1.306 & 1.690 & 2.030 & 2.133 & 2.438 & 2.724 & 3.340 & 3.591 \\
    & \textbf{40} & 0.126 & 0.255 & 0.388 & 0.529 & 0.681 & 0.851 & 1.050 & 1.303 & 1.684 & 2.021 & 2.123 & 2.423 & 2.704 & 3.307 & 3.551 \\
    & \textbf{50} & 0.126 & 0.255 & 0.388 & 0.528 & 0.679 & 0.849 & 1.047 & 1.299 & 1.676 & 2.009 & 2.109 & 2.403 & 2.678 & 3.261 & 3.496 \\
    & \textbf{60} & 0.126 & 0.254 & 0.387 & 0.527 & 0.679 & 0.848 & 1.045 & 1.296 & 1.671 & 2.000 & 2.099 & 2.390 & 2.660 & 3.232 & 3.460 \\
    & \textbf{120} & 0.126 & 0.254 & 0.386 & 0.526 & 0.677 & 0.845 & 1.041 & 1.289 & 1.658 & 1.980 & 2.076 & 2.358 & 2.617 & 3.160 & 3.373 \\
    & \textbf{\(\infty\)} & 0.126 & 0.253 & 0.385 & 0.524 & 0.675 & 0.842 & 1.036 & 1.282 & 1.645 & 1.960 & 2.054 & 2.327 & 2.576 & 3.091 & 3.291 \\
    \hline
\end{tabular}
}

\end{document}

#Como compilar e gerar uma imagem svg (Scalable Vector Graphics) para colocar em html
#pdflatex EsbocoTikz.tex
#pdf2svg EsbocoTikz.pdf EsbocoTikz.svg

    
\end{document}
%*************** FINAL DO DOCUMENTO ***************
